\documentclass{article}
\usepackage[utf8]{inputenc}
\usepackage{amsmath,amsthm,amssymb}

\begin{document}

\textit{Task 3.3}
In the NAM model for class 10 wave 2, the network autocorrelation parameter ($\rho_1$) was estimated at $0.613$ and is highly statistically significant ($p < 0.001$). Thus, the data strongly supports the hypothesis that a student's ending literacy score is positively associated with the literacy scores of their friends. 

Furthermore, when controlling for these network effects, the exogenouous variables gender ($p = 0.199$) and socioeconomic status/HISEI ($p = 0.902$) were not statistically significant predictors of a student's ending literacy score.

\textit{Task 3.4}  With the comparison between MR-QAP and NAM, we can observe the difference between social ties and individual outcomes.
In MR-QAP, gender homophily was highly significant predictor for friendship ties ($p < 0.001$). Thus, students 
tended to form same-gender groups. However, in NAM, gender was not statistically significant for predicting the student's ending literacy score.

When it comes to socioeconomic status, both NAM and MR-QAP didn't display statistical significance: HISEI did not affect a student's consideration for 
friendship ($p = 0.658$), nor did it predict high ending literacy score in NAM ($p = 0.902$).

\textit{Task 3.6}
Analysing the network autocorrelation parameter, we see that classrooms 2, 3, 4, 8, 10 and 11 have positive estimates, implying that overall, friendships are corellated with similar literacy scores. However, this cannot be generalised since only classrooms 3, 10 and 11 have storng statistical significance. For the rest of the clasrooms, the network effect is not strong enough to be discerned from the random noise. While there is a positive trend, the statistical significance of the autocorrelation hypothesis is not generalizable across all classrooms.

\end{document}
